\chapter{Landscape chart atlases for printing and presentations}\label{ch:atlases}

The companion files \path{qPCR_control_charts_realscale_A4_landscape.pdf} and
\path{qPCR_control_charts_realscale_A3_landscape.pdf} contain the complete application-exported chart set,
one graph per page. The A4 atlas is sized for landscape printing; the A3 atlas uses the same page layout at a larger
physical size for presentation and poster use. Both include every exported instrument and chart variant, including
companion subgroup-spread charts. The figures and their data are synthetic.

\section{How to read a point against the limits}
Each chart labels its measurement and unit, instrument, series, centre line and any displayed lower and upper control
limits. The short interpretation below the graph describes what the metric represents. Red points indicate values
outside a displayed limit. Amber points indicate a pattern rule inside the limits where that rule is enabled. These
are prompts for review, not automatic diagnoses.

\begin{infobox}[title={Below minimum or above maximum}]
\textbf{Below the lower limit:} the value is unusually low relative to the historical series. For a low-is-better measure,
this may be a favourable shift, but an unexpected change still needs context. For a high-is-better measure, a low value
is in the adverse direction. For a target-centred measure, either direction may be meaningful.

\textbf{Above the upper limit:} the value is unusually high relative to the historical series. For a high-is-better
measure, it may be favourable if the change is expected; for a low-is-better measure, it is in the adverse direction.
Target-centred measures can warrant review on either side.

\textbf{A control limit is not automatically a minimum or maximum allowed by the method.} It estimates variation in a
particular historical series. Where the approved SOP or vendor specification supplies acceptance limits, apply those
separately. Check the instrument, program, reagent lot, extraction batch, operator and maintenance record before drawing
a conclusion. A chart without enough comparable history may have no reliable limits yet.
\end{infobox}

The atlas footer reports whether values cross the displayed statistical limits and adds the catalogue's metric-specific
interpretation. A configured minimum or maximum from the chart catalogue is shown separately when one exists. Limits are
instrument-specific; do not compare one instrument's centre line with another's or combine different programs without
documented comparability.

\section{Selecting graphs for a presentation}
The A3 PDF is already in landscape and keeps each chart's title, unit, interpretation and limit explanation on the same
page. Use one chart per slide when the points, series and annotations need to remain readable. For an overview slide,
choose a small set that answers one question: thermal performance (heating, cooling, settling and transition time),
optics (reference, signal and exposure), process integrity (control and edit findings), or operator/method behaviour
(replicate precision and negative-control rate). State that the data are synthetic and that control limits are
historical process limits, not validated acceptance criteria.

The source chart bundle is \path{generator/out/DEMO_control_charts_realscale.zip}; the offline script
\path{generator/make_chart_atlases.py} rebuilds both page sizes from the application's SVG/CSV export.


